\documentclass[11pt]{article}
\usepackage[T1]{fontenc}
\usepackage{lmodern,amsmath,amssymb,amsthm,geometry,microtype,xcolor}
\usepackage[colorlinks=true,urlcolor=blue!55!black,linkcolor=black]{hyperref}
\geometry{margin=.93in}
\setlength{\emergencystretch}{2em}
\newtheorem{theorem}{Theorem}
\newtheorem{lemma}{Lemma}
\newcommand{\dd}{\,dl}
\newcommand{\dt}{\,dt}
\newcommand{\Rr}{\mathcal R}
\newcommand{\Ll}{\mathcal L}
\newcommand{\E}{\mathbb E}
\newcommand{\eps}{\varepsilon}
\newcommand{\norm}[1]{\lVert #1\rVert}
\newcommand{\repo}{https://github.com/TodManlaibaatar/OOD-reversals-theory/blob/12aa3b6f48f83a097600bcf9b44d93399e2f9697}
\newcommand{\src}[2]{\href{\repo/09_canonical_regime_AN04/AN04_#1_v1.tex}{#2}}
\title{AN04: From the first small weak response to the central E5 depth comparison}
\author{Version 2 \quad|\quad Research draft for review}
\date{5 November 2026}
\begin{document}
\maketitle
\begin{abstract}
At the central critical initialization, a fixed small weak-response hit has bounded total mass.
Dissipation there gives a fixed loss gap and places the first measured half response within bounded time.
We present the argument for carrying the signed strong-minus-weak history of exposed fourth-quadrant labels to the coherent witness time. The retained Q4 contribution absorbs the adverse vertical output and then dominates the horizontal cost. The conclusion concerns one endpoint of the initialized flow. The companion review gives the hypotheses, source labels and order of choices for each lemma.
\end{abstract}
\section{Claim, clocks and source results}\label{an04ce5v2:setup}
Use normalized time and the equally weighted mixture of $P_1=N(e_1,q^2I)$ and $P_2=N(\rho e_2,q^2I)$ with identity target $D$. Write $\E_j$ for expectation on $P_j$ and $\E$ for mixture expectation.
Fix $\rho=2/3$, $\zeta=5\pi/9$, $\lambda=\rho^2$, $L=\lambda/2$, $\omega=(1-\lambda)/2$ and $\ell=\log(1/q)$. With $c>0$ fixed take $S_0=cq^{18/5}$. The initial vectors are $U_\alpha=W_\alpha=\sqrt{S_0}(\cos\alpha,\sin\alpha)$, with $dl=d\alpha/(2\pi)$. The quadrant and half-circle names always refer to this initial label. With $a_+=\max(a,0)$, put
\[
 \begin{gathered} f(D)=\int W(U^\top D)_+\dd,\qquad r=D-f,\qquad
 S=\int |U|^2\dd=\int |W|^2\dd,\\
 \Ll_j=\tfrac12\E_j|r|^2,\quad \Ll=(\Ll_1+\Ll_2)/2.\end{gathered}
\]
Balance holds label by label. For $\Theta=(U,W)$ the exact characteristics are
$U'=\mathsf R^\top W$, $W'=\mathsf R U$ with
$\mathsf R=\E[rD^\top{\bf1}_{U^\top D>0}]$ for the receiving label.
Thus $\Ll'=-\|\Theta'\|^2$ and $S'=2\E[f\cdot r]\le \E|D|^2/2$.
All unlabeled norms of parameters are in $L^2(dl)$. Set $Q_1=1+q^2$, $Q_2=Q=\lambda+q^2$ and
\[
 b_j=\E_j[D_j(U^\top D)_+]/Q_j,\quad v_j=\nabla_U b_j,\quad
 \theta_j=(v_j)_j,\qquad
 F=\begin{pmatrix}m&h\\z&n\end{pmatrix},\quad Fe_j=\int W b_j\dd.
\]
Let $t_h$ be the first half mass of right ancestry and $t_a=\inf\{t\ge t_h:n(t)=a\}$ for a weak response level $a$. Thus $t_{1/2}$ is measured from the actual flow. The central coherent witness has offset $s_*$, comparison interval $J$, and $e_*=\min_{s\in J}\mathcal E_{\rm coh}(s)=\mathcal E_{\rm coh}(s_*)<0$. Define $\xi=(\sin(q\zeta),-\cos(q\zeta))$ and $T=t_{1/2}+s_*$.
\begin{theorem}[Central initialized endpoint]\label{an04ce5v2:main}
For every fixed $c>0$ and all sufficiently small $q$, the first half response is finite. At $T$,
\[
 m\ge\tfrac12,\qquad n\ge\tfrac14,\qquad
 \Ll_j(T)/\Ll_j(0)\le169/196\quad(j=1,2).
\]
The scaled depth satisfies
$e_q(T)=(|\xi-f_T(\xi)|^2-1)/(2q^2)\to-\infty$ as $q\downarrow0$.
In particular $e_q(T)<e_*$. The unscaled probe loss tends to $1/2$.
\end{theorem}
\paragraph{Inputs from the theory repository.}
\src{trained_half_mass_bridge_donors}{HBD} and \src{bridge_return_before_weak_learning}{BRW}
give $t_h=\log(4/S_0)+o(1)$. At $t_h+b$ for each fixed $b\ge0$, the strong mass and regression tend to $M(b)=(1+e^{-b})^{-1}$, the right radial density tends to $4M(b)\cos^2\alpha$, and $4\Ll\to (1-M(b))^2+\lambda$. The Cartesian entrance envelopes have horizontal scale $H_0=\sqrt{S_0}e^{(1/2+q^2)t}$ and vertical scale $V_0=\sqrt{S_0}e^{(L+q^2)t}$. In particular both second components are $O_b(q)$ and the left radial mass per label is $O_b(q^2(1+t_h)^2+S_0)$ at fixed delay.
\src{core_output_descent}{COD} gives the resident rank ceiling and bottom-path bounds.
\src{transverse_balance_and_weak_clock}{TBW} gives the seed
\[
 q^{-2}\int_{\rm Q2}[(U_2+W_2)/2]^2\dd\ \longrightarrow\ c^{2\omega}(4e^b)^\lambda/8 .
\]
On compact Q2 interiors the seed divided by $q$ tends to $c^\omega(4e^b)^L\sin\alpha$ and both first components are $o_b(q)$. Fix $\eta$ as in TBW's small-hit proof, decreasing it if needed. On a fixed positive-length compact Q2 arc before $t_\eta$, $v=(U_2+W_2)/2>0$, $|U_1|,|W_1|,|U_2-W_2|/2\le v/2$, and $v'\ge\lambda v/4$ after a sufficiently late fixed entrance. Its loss budget gives $m>99/100$ there.
\src{residual_exchange_channels}{REC} and \src{cap_free_cross_response}{CFCR} give the Gaussian regression flow and shape estimates. We use the transverse energy identity of \src{gate_transport_and_transverse_memory}{GTM}. The \href{\repo/04_witnesses_persistence/AN06_learned_orbit_witnesses_v2.tex}{WIT central margins} give $s_*>-\log(3)/\lambda$. Constants may depend on $c$, on fixed witness data and on choices already fixed. The five expanded lemmas appear in dependency order: the cone, charged energy, shrinking bands, retained labels and final continuation. All estimates concern the initialized population flow.

\section{Mass and time at the first small hit}\label{an04ce5v2:mass}
Write $t_0=t_h+b_0$ for a sufficiently late fixed entrance, $s=t-t_0$,
$\tau=t_\eta$, and $E=\tfrac12\int (U_2^2+W_2^2)\dd$. Stop at the first of
$\tau$, $S=4$, and $t_d=t_0+(2/Q)\log(d/q)$ for fixed $0<d<1$.
On this interval
\begin{equation}\label{an04ce5v2:energy}
 E'\le QE+C(q\sqrt E+q^2),\qquad |n|\le E+Cq\sqrt E .
\end{equation}
Indeed $b_2=(U_2)_++O(q|U|)$ gives the second bound by Cauchy--Schwarz.
The exact transverse identity is $E'=\E[D_2r^\top\partial_2 f+r_2 f_2]$.
On $P_1$, $\|\partial_2 f\|_\infty\le\sqrt{2SE}$, $D_2$ has norm $q$, and $r_2 f_2\le D_2^2/4$. On $P_2$, Euler's identity for the homogeneous map $f$ writes the integrand as
\[
 (D_1-f_1)(f_1-D_1\partial_1 f_1)+2r_2 f_2-D_1r_2\partial_1 f_2.
\]
The first product is at most $D_1^2(1-\partial_1 f_1)^2/4$ by completing the square. The last term costs $O(q\sqrt E)$, and the middle term has expectation at most $2Qn$. The cluster weight $1/2$ gives \eqref{an04ce5v2:energy}. Dividing by $2\sqrt{E+q^2}$ and integrating yields
\begin{equation}\label{an04ce5v2:envelope}
 \sqrt E\le C_{b_0}q e^{Qs/2},\quad
 \int_{t_0}^t\sqrt E\le C_{b_0}d,\quad \int_{t_0}^t E\le C_{b_0}d^2 .
\end{equation}
Projecting $S'$ on the principal cluster coordinates and completing the two residual squares gives $S'\le Q_1m(1-m)+Qn(1-n)+q^2/2$. Moreover
\[
 m'=a_m(1-m)+O(\sqrt E+q),\qquad
 a_m=\tfrac{Q_1}2\|\nabla m\|^2\ge Q_1m^2/S.
\]
Here $|\langle\nabla m,\nabla n\rangle|\le Cq S$ since $|v_1^\top v_2|\le Cq$, $|h|+|z|\le C(\sqrt E+q)$, and the projected shape error is $O(q)$ on the stop. Euler's identity $\langle\Theta,\nabla m\rangle=2m$ gives the lower bound. Integrating the damped equation bounds the positive deficit cost by $C(1-m(t_0))_++C\int(\sqrt E+q)$. Hence
\[
 S(t)\le S(t_0)+C(1-m(t_0))_++C_{b_0}(d+d^2)+o_q(1)<4 .
\]
Choose $b_0$ first using $S(t_0),m(t_0)\to M(b_0)$, then $d$, then $q$. This removes the mass stop.

\begin{lemma}[Q2 cone on both sides of the amplitude cutoff]\label{an04ce5v2:conelemma}
With $b_0$ sufficiently large, $\eta>0$ sufficiently small and fixed, and then $d>0$ fixed as above, the initialized first small hit satisfies
\begin{equation}\label{an04ce5v2:smallhit}
 \tau-t_h=\ell/L+O(1),\qquad
 \sup_{[t_0,\tau]}S+\int_{t_0}^{\tau}(|1-m|+|n|)\dt=O(1).
\end{equation}
On each fixed compact Q2 interior the positive cone below holds up to $\tau$ for small $q$.
\end{lemma}
\begin{proof}
Choose $b_0$ also large enough, and then $\eta<1/8$ small enough, that TBW \texttt{lossbudget} has square root less than $\lambda/400$ uniformly to the first $\eta$ hit. Fixed-delay BRW puts every earlier post-strong-half response up to $t_0$ below $\eta$ for small $q$. More explicitly, $4\Ll(t_0)\to(1-M(b_0))^2+\lambda$, and $\E_2 r_2^2\ge Q(1-n)^2\ge Q(1-\eta)^2$ on the no-hit interval. Thus
$\|r\|_{P_1},\|r_1\|_{P_2}\le\lambda/400$; $\|r_2\|_{P_2}\le C$ by initial loss. The regression identity then gives $m>99/100$. These choices can be made jointly with the preceding strict mass bound: its damping constants on $S\le4$, $m>99/100$ are absolute, while $d$ is chosen only after the fixed entrance.

Here is the budget for each field entry, including after $t_d$. Write $r_{ip}=\|r_i\|_{P_p}$ and $\Gamma={\bf1}_{U^\top D>0}$. The all-gate Cauchy--Schwarz bounds for the first three entries are
\[
 \begin{aligned}
 2|\mathsf R_{11}|&\le\sqrt{Q_1}\,r_{11}+q r_{12},\\
 2|\mathsf R_{12}|&\le q r_{11}+\sqrt Q\,r_{12},\\
 2|\mathsf R_{21}|&\le\sqrt{Q_1}\,r_{21}+q r_{22}.
 \end{aligned}
\]
Each is sufficient for $|\mathsf R_{ij}|\le\lambda/100$ for small $q$ (only $r_{22}$ uses the full loss bound). When $U_2/|U|\ge1/\sqrt2$, $\Pr_2(\Gamma=0)\le e^{-\rho^2/(4q^2)}$. The missing $P_2$ correlation costs at most $r_{22}(\E_2 D_2^4)^{1/4}\Pr_2(\Gamma=0)^{1/4}$. The ungated correlation is exactly $Q(1-n)$; the entire $P_1$ gate term costs at most $q r_{21}$. Hence
\[
 \mathsf R_{22}\ge Q(1-n)/2-Cq-Ce^{-c'/q^2}\ge\lambda/3
\]
on that cone through the full no-hit interval. These four estimates use loss only.

Let $v=(U_2+W_2)/2$, $d_v=(U_2-W_2)/2$. TBW's fixed-delay entrance on each chosen compact has $v/q$ bounded below by a positive constant, $U_1,W_1=o(q)$, $d_v=o(q)$. Thus $v>0$ and $|U_1|,|W_1|,|d_v|<v/2$ at $t_0$. On the closed faces $U_2,W_2\in[v/2,3v/2]$ and $U_2/|U|\ge1/\sqrt2$. The exact equations imply
\[
 v'=\mathsf R_{22}v+(\mathsf R_{12}W_1+\mathsf R_{21}U_1)/2
 \ge(\lambda/3-\lambda/200)v>\lambda v/4 .
\]
At $U_1=\pm v/2$, $D^+|U_1|\le |\mathsf R_{11}W_1|+|\mathsf R_{21}W_2|\le(2\lambda/100)v<v'/2$. At $W_1=\pm v/2$ use $|\mathsf R_{11}U_1|+|\mathsf R_{12}U_2|$ for exactly the same bound. At $d_v=\pm v/2$, its derivative is $-\mathsf R_{22}d_v+(\mathsf R_{12}W_1-\mathsf R_{21}U_1)/2$, so $D^+|d_v|\le(-\lambda/6+\lambda/200)v<v'/2$. These strict inequalities hold at intersections of any of the six faces. $v\ge v(t_0)e^{\lambda(t-t_0)/4}$ excludes the zero face. Each label therefore stays in the cone until the first hit, inclusively by continuity. In particular an arrival at $t_d$ with $\tau>t_d$ is an interior continuation time for the same loss argument; a simultaneous hit ends it there.

To obtain the sharper clock use the bounded-mass estimates \emph{only} up to $\min(t_d,\tau)$. REC's shape decomposition in the receiving field gives
$\mathsf R_{12}=\tfrac12Q_1(1-m)(v_1)_2-\tfrac12 Q h(v_2)_2+O(q)$ and
$\mathsf R_{21}=-\tfrac12Q_1 z(v_1)_1+\tfrac12Q(1-n)(v_2)_1+O(q)$.
Each remainder is bounded by $\frac12\sum_p\|e^{(p)}\|_{P_p}\|D_j\|_{P_p}=O(q)$ on $S\le4$. Thus $|\mathsf R_{12}|+|\mathsf R_{21}|\le C(\sqrt E+q)$. Combining the cone faces, the sharper lower bound $Q(1-n)/2-Cq$ and \eqref{an04ce5v2:energy} yields
$v'\ge[Q/2-C(E+\sqrt E+q)]v$ on this restricted interval. Integration using \eqref{an04ce5v2:envelope} gives a lower multiplier $c_{b_0,d}e^{Q(t-t_0)/2}$ since $q(t_d-t_0)=o(1)$. At $t_d$ this exponential equals $d/q$. Therefore, if $\tau>t_d$, a \emph{fixed positive-measure} interior arc $J_0$ has $v(t_d)\ge k_{b_0,d,J_0}>0$. For $0\le u\le\tau-t_d$ the already checked loss-budget cone and $v^2\le |U|^2$ by balance give
\[
 |J_0|k_{b_0,d,J_0}^2 e^{\lambda u/2}
 \le\int_{J_0}v(t_d+u)^2\dd\le S(t_d+u)\le4+Cu .
\]
Here $|J_0|=\int_{J_0}dl$, and the last inequality uses the global $S'$ bound. The exponential exceeds the line at a fixed finite $u$, so the hit occurs within $O(1)$ after $t_d$. This argument needs no sharpened off-diagonal estimate beyond $t_d$.

If $\tau\le t_d$, \eqref{an04ce5v2:energy} and $n(\tau)=\eta$ give $\sqrt{E(\tau)}\ge c_\eta>0$ for small $q$, so \eqref{an04ce5v2:envelope} gives $\tau-t_0\ge(2/Q)\ell-O(1)$. The upper bound is $t_d$ in this case. If $\tau>t_d$, both bounds follow from $t_d-t_0=(2/Q)\log(d/q)$ and the bounded continuation. Because $q^2\ell=o(1)$, $(2/Q)\ell=\ell/L+o(1)$. Up to $\min(\tau,t_d)$, the damped $m$ equation integrated in absolute value gives $\int |1-m|\le C(|1-m(t_0)|+\int(\sqrt E+q))$, and \eqref{an04ce5v2:energy} bounds $\int |n|$. The possible remaining interval is bounded and both regressions are bounded by loss. This proves \eqref{an04ce5v2:smallhit}. A fixed further continuation has bounded mass by the same linear inequality.
\end{proof}


\section{A fixed loss gap and the first half response}\label{an04ce5v2:half}
Let $A=\operatorname{diag}(\sqrt{Q_1},\sqrt Q)$, $u=A(b_1,b_2)^\top$, and $D_{\Ll}=4\Ll$.
Cluster projection gives $D_{\Ll}=\|(I-F)A\|_F^2+\chi$ with $0\le\chi\le Cq^2(1+S)^2$. Indeed $r$ is $(1+S)$-Lipschitz on each cluster, and approximation by its mean-center value times $D_j/\E_jD_j$ leaves $O(q(1+S))$ in $L^2(P_j)$. Since $FA=\int Wu^\top\dd$, $\int uu^\top\dd\succeq (FA)^\top(FA)/S$. The output gradient is $\dot W=(I-F)Au/2+N_W$ with $\|N_W\|\le Cq\sqrt S(1+S)$. Consequently
\begin{equation}\label{an04ce5v2:gram}
 \|\dot W\|\ge\frac{\sigma_{\min}(FA)}{2\sqrt S}\|(I-F)A\|_F-Cq\sqrt S(1+S).
\end{equation}
The bounded mass in \eqref{an04ce5v2:smallhit} bounds $|n'|$ uniformly. There is a fixed $a>0$ with $\eta/2\le n\le\eta$ on $[\tau-a,\tau]$. For every sufficiently late fixed $b$ the entrance at $t_h+b$ precedes this interval. Loss monotonicity and then $b\to\infty$ give $\limsup_{q\to0}(D_{\Ll}(\tau-a)-Q)\le0$. With $\gamma=\sqrt{Q/Q_1}$, the exact square completion reads
\[
 2\gamma\det F=(Q-D_{\Ll}+\chi)/Q_1+
 (m-1+\gamma n)^2+(z-\gamma h)^2+2\gamma(1-\gamma)n .
\]
Thus $\det F$ is bounded below on a fixed time interval. Since $\|FA\|$ is bounded there, its least singular value is bounded below too. The error norm in \eqref{an04ce5v2:gram} is at least $\sqrt Q(1-\eta)$, so $D_{\Ll}'\le-c_0<0$ there. In particular
\begin{equation}\label{an04ce5v2:gap}
 D_{\Ll}(\tau)\le Q-\Delta
\end{equation}
for a fixed $\Delta>0$. Afterwards $\sigma_{\min}(FA)\ge\sqrt Q-\sqrt{Q-\Delta}$ by the triangle inequality. Before half response the remaining error norm is at least $\sqrt Q/2$. On $[\tau,\tau+K]$, $S\le B+C(t-\tau)$. Choose $K$ so that the dissipation lower bound proportional to $\int_0^K (B+Cu)^{-1}du$ exceeds all available loss, then take $q$ small enough to absorb the error in \eqref{an04ce5v2:gram} throughout that fixed interval. This forces
\begin{equation}\label{an04ce5v2:halfclock}
 0\le t_{1/2}-\tau\le K,\qquad t_{1/2}=t_h+\ell/L+O(1).
\end{equation}
The argument requires neither monotonicity of $n$ nor a small cross-response cap.

\section{Axis separation and the learning guards}\label{an04ce5v2:guards}
\begin{lemma}[Charged energy and iterated axis separation]\label{an04ce5v2:axislemma}
Fix a finite $K$, $t_f=t_{1/2}+K$, and a closed interior arc $J_2\Subset{\rm Q2}$.
Charge coordinate $i=1$ on $D_1=J_2$ and coordinate $i=2$ on $D_2=J_2^c$ in both layers. In this lemma $D_i$ denotes a label set when used as an integration domain. Let $P$ be this projection and $E_J=\|P\Theta\|^2/2$. There is a constant independent of $J_2$ such that on $[t_0,t_f]$
\begin{equation}\label{an04ce5v2:defect}
 E_J'\le [Q+C(|n|+|1-m|)]E_J+C(q\sqrt{E_J}+q^2).
\end{equation}
The supremum of $\sqrt{E_J}$ on $[t_h+b,t_f]$ has iterated upper limit zero on taking first $q\downarrow0$, then $b\to\infty$, then $J_2\uparrow{\rm Q2}$. The right vertical norm, $|z|+|h|$ and $|T_1|$ tend uniformly to zero on $[t_0,t_f]$, where $T_1=\int W_1W_2|v_1|^2\dd$.
\end{lemma}
\begin{proof}
Write $g_{ij}=\nabla F_{ij}=(W_i v_j,e_i b_j)$ and $j'=3-j$. The small-hit clock and fixed loss gap in Sections~\ref{an04ce5v2:mass}--\ref{an04ce5v2:half} give a common $S\le S_*$ and a common bound for $\int_{t_0}^{t_f}(|n|+|1-m|)$. Loss monotonicity bounds each regression uniformly. Gaussian homogeneity gives the exact identity $b_j=U_j(v_j)_j+U_{j'}(v_j)_{j'}$, also at $U=0$ with the zero-gate convention. REC gives $|(v_j)_{j'}|\le Cq$ and $|b_j-(U_j)_+|\le Cq|U|$. These identities retain the derivative even on a center-gate boundary.

For $i=j$, the pairing with $P\Theta$ on $D_j$ has integrand
$W_j(b_j+U_j(v_j)_j)$; on $D_{j'}$ it is $W_j U_{j'}(v_j)_{j'}$. Hence
\[
 \langle P\Theta,g_{jj}\rangle
 =2\int_{D_j}W_j(U_j)_+\dd+e_j,\qquad
 |e_j|\le Cq\left(\int_{D_j}|W_j||U|\dd+
              \int_{D_{j'}}|W_j||U_{j'}|\dd\right)\le Cq\sqrt{E_J}.
\]
The last estimate uses $S_*$ and, in each integral, one charged factor. The leading integral in absolute value is bounded by $\int_{D_j}(W_j^2+U_j^2)\dd\le2E_J$ including its coefficient two. This accounts separately for $g_{11}=g_m$ and $g_{22}=g_n$.

For $i\ne j$, on $D_i$ the integrand is $W_i(b_j+U_i(v_j)_i)$; on $D_j$ it is $W_iU_j(v_j)_j$. Subtract $F_{ij}$ and use Euler's identity on $D_j$ to get
\[
 \langle P\Theta,g_{ij}\rangle-F_{ij}
 =\int_{D_i}W_iU_i(v_j)_i\dd-\int_{D_j}W_iU_i(v_j)_i\dd,\qquad
 |\langle P\Theta,g_{ij}\rangle-F_{ij}|\le CqS.
\]
Moreover $|F_{ij}|\le C\sqrt{E_J}+Cq$: approximate $b_j$ by $(U_j)_+$ and use $W_i$ on $D_i$, $U_j$ on $D_j$ as the charged factor. This covers $(i,j)=(2,1),(1,2)$ with their own coefficients. Their flow contributions are $-Q_1z^2/2-Qh^2/2+O(q\sqrt{E_J}+q^2)$. The $m$ and $n$ terms have coefficients $Q_1(1-m)/2$ and $Q(1-n)/2$; $|1-n|\le1+|n|$ gives the leading $QE_J$. Finally $|\langle P\Theta,N_{\rm sh}\rangle|\le\sqrt{2E_J}\,Cq\sqrt S(1+S)$ by REC's data-space shape projection. Adding the four pairings and this error proves \eqref{an04ce5v2:defect}.

Put $Y=(E_J+q^2)^{1/2}$. Division gives $Y'\le[Q/2+C(|n|+|1-m|)]Y+Cq$. From $t_b=t_h+b$ onwards variation of constants, $\int_0^\infty e^{-Qu/2}du=2/Q$, and the integrated budget give
\[
 \sup_{[t_b,t_f]}Y\le C_K e^{Q(t_f-t_b)/2}(Y(t_b)+q)
 \le C_K e^{-Lb}\big(1+\sqrt{E_J(t_b)}/q\big).
\]
Here $b\ge b_0$ is fixed before $q$; $t_f-t_h=\ell/L+O_K(1)$ and $q^2\ell\to0$. The last constant is independent of $b,J_2$.

We compute the entry on arbitrary subsequences at this \emph{fixed} $b$. HBD \texttt{clock} and BRW \texttt{field, delay} transfer compact right paths with mass $\mu_b=M(b)\pi$, $d\pi=4\cos^2\alpha\,dl$. The integrand of the right part of $E_J/q^2$ on a compact right interior is the physical radial mass times
$[(\sin(qx)/q)^2+(\sin(qy)/q)^2]/2$. It converges in averaged label law to $(x^2+y^2)/2$ integrated against $\mu_b$. At this same $b$ the all-label vertical bound $C_b q$ charges a strip of length $\sigma$ at most $C_b^2\sigma$ in scaled energy. Take its limsup at fixed $\sigma,b$, then remove this auxiliary strip at fixed $b$. The identical argument removes the Q2 and Q3 endpoint strips in the next calculation. On compact Q3 the squared radius is $O_{b,\sigma}(S_0)$ by the proof of COD \texttt{trained}; $S_0/q^2\to0$. On $J_2$, TBW \texttt{smallhit} gives $U_1,W_1=o_{b,J_2}(q)$. On compact Q2 the positive second components, balance and TBW \texttt{seed} give
$U_2/q,W_2/q\to c^\omega(4e^b)^L\sin\alpha$. Thus the Q2 contribution outside $J_2$ divided by $q^2$ tends to $c^{2\omega}(4e^b)^{2L}\int_{{\rm Q2}\setminus J_2}\sin^2\alpha\dd$.

It remains to bound the right moment after discount, uniformly over resident extractions. COD \texttt{ceiling} bounds $e^{-2Lb}(x_+^2+y_+^2)$ by $C(1+f^{-1/3})^2 e^{-\lambda fb/8}$ at normalized upper Q1 rank $f$, with $|df|=d\pi$, and with $f=1/2$ on Q4. The envelope is integrable at zero, so its integral tends to zero. For the negative parts COD's lower face in the proof of \texttt{exposure} gives $x_-\le C\sqrt{1+b}$ by order. The bottom $y$ equation in COD \texttt{field}, its bounded entry, its damping kernel $\le1$, and BRW \texttt{budget} give $(y_a)_-\le C+C\int_0^b e^{Lu}(1+u)^{-2/3}du$. After multiplication by $e^{-Lb}$, splitting at $b/2$ bounds this by $O(e^{-Lb/2}+(1+b)^{-2/3})$. Order bounds every $y_-$ by $(y_a)_-$. These bounds hold for every extracted law so also bound the original limsup. It follows that
\[
 \limsup_{b\to\infty}\limsup_{q\downarrow0}
 \sup_{[t_b,t_f]}\sqrt{E_J}
 \le C_K\left(\int_{{\rm Q2}\setminus J_2}\sin^2\alpha\dd\right)^{1/2}.
\]
Exhaust Q2 only now. On this late interval the right vertical norm is at most $2\sqrt{E_J}$; the two cross responses have the bound derived above. For $T_1$ use $W_1$ on $J_2$, $W_2$ on its complement and $|v_1|\le C$. On $[t_0,t_b]$, for each fixed $b$ the all-label vertical envelope is $C_bq$, the mass bounded, $b_2=(U_2)_++O(q|U|)$ and $\|b_1\|\le C$. Thus the same quantities are $O_b(q)$ there. Given any tolerance choose the arc and fixed delay according to the iterated bound, then take $q$ small on both intervals. This proves uniform convergence on their union.
\end{proof}


For the strong response, the scalar four-regression equation has coefficient $a_m=Q_1\|\nabla m\|^2/2\ge Q_1m^2/S$. It is bounded below on the whole interval: before $\tau$ use Lemma~\ref{an04ce5v2:conelemma}, afterwards $Q_1(1-m)^2\le Q-\Delta$. The forcing in $(1-m)'+a_m(1-m)$ is $O(|z|+|h|+q)$ since $\langle g_m,g_n\rangle=O(qS)$ and the other gradients are bounded. It is uniformly $o(1)$ by Lemma~\ref{an04ce5v2:axislemma}. The entrance homogeneous term decays exponentially over the diverging time to every fixed half-centered window. Thus $m\to1$ there.

For the weak gradient fix such a window, with $K$ enlarged to include its endpoints and the bounded post-$\tau$ continuation. Put $\epsilon_J=\sup_{[t_b,t_f]}\sqrt{E_J}$. On $J_2$ up to $\tau$, the cone gives $U_2,W_2>0$ and balance gives
$|W_2-U_2|^2\le |W_2^2-U_2^2|\le W_1^2+U_1^2$.
After $\tau$, integrate the equations for $(U_2-W_2)/2$ and $(U_2+W_2)/2$ with coefficients $\mp\mathsf R_{22}$ and horizontal forcing bounded by $C(|W_1|+|U_1|)$. On this bounded interval both propagators are positive and bounded by a common constant. The difference mode in label norm, and the negative part of the sum mode, are therefore bounded by $C_K\epsilon_J$ on $J_2$, using the nonnegative sum at $\tau$. This bounds $\|\mathbf1_{J_2}(W_2-U_2)\|$ and $\|\mathbf1_{J_2}(U_2)_-\|$ there. On $J_2^c$ both vertical norms cost at most $2\epsilon_J$. Take the iterated limits of the lemma to obtain $\|W_2-(U_2)_+\|=o(1)$ on the centered window. Consequently $\|W_2-b_2\|=o(1)$. Also $|b_2(1-|v_2|^2)|\le Cq|U|$: use $|b_2-(U_2)_+|\le Cq|U|$, $|(v_2)_1|\le Cq$, and for positive $U_2$ bound $1-\theta_2$ by the weighted Gaussian wrong-gate tail $Ce^{-c(U_2/(q|U|))^2}$. Hence
$\int b_2^2|v_2|^2\dd=\int b_2^2\dd+O(qS)$. In the exact formula
$\|\nabla n\|^2=\int(b_2^2+W_2^2|v_2|^2)\dd$ replace $W_2$ by $b_2$ using the bounded norms, and use $n=\int W_2 b_2\dd$. This gives $\|\nabla n\|^2=2n+o(1)$. In CFCR's exact flow the $g_m$ pairing costs $O(qS)$, the $g_z,g_h$ terms vanish with $z,h$, and the shape remainder costs $O(q)$. Thus $n'=Q n(1-n)+o(1)$ uniformly on the window. Using the measured centering $n(t_{1/2})=1/2$ gives uniformly for bounded $s$
\begin{equation}\label{an04ce5v2:logistic}
 (m,z,h,n)(t_{1/2}+s)\longrightarrow(1,0,0,N(s)),\qquad
 N(s)=(1+e^{-\lambda s})^{-1}.
\end{equation}
This also bounds $1-n$ away from zero through $T$ by combining the fixed continuation and the pre-hit bound. At $s_*$, $N(s_*)>1/4$. Cluster projection and $f_0(D)=S_0D/4$ give loss ratios tending to $0$ and $(1-N(s_*))^2<9/16<169/196$. This proves both response and both loss guards.

\section{Probe exposure and the signed history to be estimated}\label{an04ce5v2:history}
Set $p_\alpha=(U_\alpha(T)^\top\xi)_+$, $F_i=f_{T,i}(\xi)$, $H=F_1/q$ and $w=Q(1-n)$. Before $t_d$ on the stopped interval, \eqref{an04ce5v2:envelope} bounds $\int\sqrt E\dt$. After that interval ends there are at most $O(1)$ units left to $T$, by the clock comparison. Consequently
\[
 \int_{t_0}^T(\|\Theta'\|+\|r\|_{P_1})\dt\le C,\qquad
 \|\mathcal U\|\le C,\quad
 \mathcal U_\alpha=\sup_{[t_0,T]}(|U_\alpha|+|W_\alpha|).
\]
Here the regression terms in $\|\Theta'\|$ are bounded by $C(|1-m|+|h|+|z|+\sqrt E+q)$ since $\|\nabla n\|\le C(\sqrt E+q)$. The strong residual costs $C(|1-m|+|z|+q)$ by shape projection. Minkowski's inequality bounds the time supremum in label norm.
On an excursion below $U_2=-A_0q\sqrt\ell\,\mathcal U_\alpha$, the $P_2$ gate has Gaussian upper tail and costs at most $Cq^2\mathcal U_\alpha$ per time in $U_2'$, for a sufficiently large fixed $A_0$. The $P_1$ cost is at most $Cq\mathcal U_\alpha\|r\|_{P_1}$ by Cauchy--Schwarz. Integrate from the last crossing or $t_0$, using the $O_{b_0}(q)$ entrance. Uniformly up to $T$,
\[
 (U_2)_-\le C_{b_0}q+Cq\sqrt\ell\,\mathcal U_\alpha,\qquad
 \|p\|\le Cq\sqrt\ell,\qquad |F_1|+|F_2|\le Cq\sqrt\ell .
\]
Define the strong-minus-weak history and the signed budget by
\begin{align}
 G_\alpha&=\int_0^T (Q_1 z b_1-w b_2)\dt,\nonumber\\
 M_q&=\frac{\int_{\rm Q4}pG_\alpha\dd-\int_{{\rm Q4}^c}p(-G_\alpha)_+\dd}{2q^2}.
 \label{an04ce5v2:budget}
\end{align}
The sign follows from $W_2'=(w b_2-Q_1zb_1)/2+\varepsilon_\alpha$. Orthogonal projection of both the residual and feature off $D_j$ in each cluster bounds $|\varepsilon_\alpha|\le Cq^2(1+S)|U_\alpha|$. Bounded mass through $T=O(\ell)$ therefore gives
$\int p|W_2(T)+G_\alpha/2|\dd\le C\|p\|(\sqrt{S_0}+q^2\ell)=o(q^2)$.
In particular $-F_2/q^2\ge M_q-o(1)$. The exact scaled loss is
\[
 e_q=H^2/2-\frac{\sin(q\zeta)}q H+
 \left(\cos(q\zeta)+F_2/2\right)F_2/q^2 .
\]
The signed endpoint criterion is
\begin{equation}\label{an04ce5v2:criterion}
 \liminf_{q\downarrow0}\left[M_q-H^2/2+\frac{\sin(q\zeta)}qH+e_*\right]>0 .
\end{equation}
For this implication, $A_q=H^2/2-q^{-1}\sin(q\zeta)H-e_*=O(1+\ell)$ and the vertical coefficient is $1+O(q\sqrt\ell)$. On the branch $F_2=o(1)$ the threshold for $-F_2/q^2$ is $A_q+o(1)$, since $q\sqrt\ell(1+\ell)=o(1)$. The estimates below will give a diverging margin in \eqref{an04ce5v2:criterion}.

\paragraph{Why the signed history is needed.}
There is also an exact pairwise diagnostic:
\[
 zF_1-mF_2=-\tfrac12\iint\det(W_a,W_b)
   (b_{1,a}p_b-b_{1,b}p_a)\,dl_a\,dl_b .
\]
It follows by antisymmetrizing the two products of integrals. To see the quantitative issue, prescribe two equal strong masses with scaled input angles $\zeta-\epsilon,\zeta+\epsilon$ and output angles $-1,+1$, respectively, and a vertical weak population. In the limiting probe formula only the first strong atom is exposed. The horizontal output divided by $q$ is $\epsilon/2$ and the vertical output divided by $q^2$ is $-\epsilon/2$. The scaled error is $\epsilon^2/8-(\zeta+1)\epsilon/2$, which approaches zero as $\epsilon\downarrow0$. A favorable pairwise sign can therefore coexist with an error above $e_*$. The initialized proof needs \eqref{an04ce5v2:criterion} with its magnitude and its charge for adverse labels.

For the remainder set $A_R=\{\cos\alpha>0\}$, $A_L=\{\cos\alpha<0\}$,
\[
 C(t)=\int (U_1)_+(U_2)_+\dd,\quad I(t)=\int_{t_0}^t C(u)\,du,\quad
 B(t)=q(1+t-t_0),\qquad t_R=t_h+\ell/L-R .
\]
Here $C(t)$ denotes the overlap and unindexed constants $C$ in inequalities denote bounds. For all sufficiently large fixed $R$, $t_R$ precedes both $t_d$ and $\tau$ by the two-sided clock estimate. Equation~\eqref{an04ce5v2:envelope} then gives $\|U_2\|+\|W_2\|+\|b_2\|\le r(t)=C_0q e^{Q(t-t_0)/2}$ up to $t_R$, with $r(t_R)\le C_1e^{-LR}$ independently of fixed $R$ for small $q$.

\section{Uniform prefix estimates}\label{an04ce5v2:prefixsec}
The following details make the comparisons quantitative before the
terminal clock. Set $d_i=W_i-b_i$ and $w=Q(1-n)$. In the characteristic
equations $U'=\mathsf R^\top W,\ W'=\mathsf R U$, the regression
decomposition, including its uniform per-receiver shape error, gives
\[
 \mathsf R=\tfrac12 Q_1(1-m,-z)^\top v_1^\top+
            \tfrac12 Q(-h,1-n)^\top v_2^\top+O(q).
\]
Here the pairs denote column vectors in the outer products. The error
follows by data Cauchy--Schwarz against $D$ times the receiver gate and
the $O(q(1+S))$ cluster shape residual. Equivalently the output-layer
equation uses $b_i$ with error $O(q^2|U|)$ by orthogonality also of the
feature remainder. We will often only need $O(q|U|)$.
In particular $\mathsf R_{22}\ge-Cq$ and
$|\mathsf R_{12}|+|\mathsf R_{21}|\le C(|z|+|h|+q)$ uniformly.

For $x=1-m$, the scalar damped equations from the four-regression flow are
\[
 x'=-a_m x+O(|T_1 z|+|h|C+q),\qquad
 z'=-a_z z+\tfrac12 w C+\tfrac12 Q_1T_1 x+O(q).
\]
Here $a_z=Q_1\|\nabla z\|^2/2\ge Q_1\|b_1\|^2/2\ge Q_1m^2/(2S)$ by $m=\langle W_1,b_1\rangle$. Both coefficients have positive upper and lower bounds. On the incoming
interval $a_m\ge Q_1m^2/S>L+\eps$ for a fixed $\eps>0$.
The integrated cost of $T_1$ times the homogeneous entrance deficit
is $O_{b_0}(q)$: $|T_1|\le Cr(t)$ on the incoming interval, and the final
fixed interval starts after the entrance deficit has decayed to $o(q)$.
The coefficient comparison can use one \emph{fixed} incoming cutoff
before $R$ is chosen. Integrating the $z$ propagator up to an arbitrary
$t$ gives weight comparable to $1\wedge(t-u)$. Duhamel in the $x$
equation can be made explicit. Write $x=x_0 e^{-\int_{t_0}^t a_m}+x_f$, $x_0=x(t_0)$, and let $\beta=\sup(|T_1|+|h|)=o(1)$. Integrating the decaying kernel in $x_f$ gives
$\int|x_f|\le C\{\beta\int |z|+\beta I+q(t-t_0)\}$ on each prefix. Hence $\int|T_1 x|\le CB+o(1)(I+\int|z|)$. In the $z$ equation the $wC$ source is nonnegative, $w$ is bounded above and below, and $z(t_0)=O(q)$. Integrate first in the receiving time; the absolute forcing error is bounded by $C\int|T_1x|+CB$ in this integral norm. Absorb the $o(1)\int |z|$ into the absolute estimate, then substitute into the negative-part and signed estimates. Thus
\begin{align}
 \int_{t_0}^t |z|&\le C(I+B),&
 \int_{t_0}^t z_-&\le o_q(1) I+CB, \nonumber\\
 \int_{t_0}^t z&\ge c\int_{t_0}^t[1\wedge(t-u)]C(u)\,du-o_q(1)I-CB.
 \label{an04ce5v2:prefix}
\end{align}
Constants here and in the next estimates are independent of the compact
right-label truncation and the large incoming offset $R$.

Differentiating $b_i$ uses $v_i$, with $|v_1^\top v_2|\le Cq$,
$b_i=(U_i)_++O(q|U|)$ and
$|b_i(1-|v_i|^2)|\le Cq|U|$. The last bound uses
$\sup_{x>0}x\Phi(-x)<\infty$ in the positive coordinate.
The $d_1$ equation contains the terms
\[
 2d_1'=-Q_1x|v_1|^2d_1-Qb_2\langle b_2,d_1\rangle+
 Q_1 z W_2 |v_1|^2-Qb_2\langle b_1,b_2\rangle+O_{L^2}(q).
\]
At $t_0$, $\norm{d_1}=O_{b_0}(q)$: integrate the same inequality from
initialization, with $\norm{d_1(0)}=O(\sqrt{S_0})$, the HBD bounds for
the cross regressions and overlap $O(H_0V_0+qH_0^2)$, and the integrated
early shape cost $O(q\int H_0)=O_{b_0}(q)$, with the entrance envelopes of Section~\ref{an04ce5v2:setup}. The growth from $x_-$ is bounded by the early mass integral.
Dropping the dissipative rank-one term and using $\int |x|\le C$ gives
\[
 \norm{d_1(t)}\le C(I+B),\qquad
 \norm{d_1(t)}\le C(r(t)I+B)\quad(t\le t_R).
\]
Since $h=\langle b_1,b_2\rangle+\langle d_1,b_2\rangle$,
$\langle b_1,b_2\rangle=C+O(q)$ and $\int\norm{b_2}\le C$, we get
\begin{equation}
 \int_{t_0}^t|h|\le C(I+B),\qquad
 h_-(t)\le C r(t)^2 I(t)+C B(t)\quad(t\le t_R).
 \label{an04ce5v2:earlyh}
\end{equation}
Similarly $2d_2'=-w|v_2|^2d_2-Q_1zb_1+QhW_1|v_2|^2+O_{L^2}(q)$, so
$\norm{d_2}\le C(I+B)$. This dissipative bound will also control costly
returns outside the retained bulk.

For clarity, the terminal weight in \eqref{an04ce5v2:prefix} can be removed
up to constants and $O(B)$. Differentiation of
$C=\int(U_1)_+(U_2)_+\dd$ a.e. gives $|C'|\le C_2(C+I+B)$:
replace $W_i$ in each diagonal characteristic term by $b_i+d_i$,
use the two norm bounds, and use
$|z|+|h|\le C_2(C+I+B)$ for the off-diagonals. For this pointwise bound use the $z$ propagator bounded by one, $\int C=I$ and the just-established prefix bound for $\int|T_1 x|$. For $h$ use $h=\langle b_1,b_2\rangle+\langle d_1,b_2\rangle$.
With $I'=C$, Gronwall on an interval of length at most two,
starting at a point of $[t-2,t-1]$ where $C\le I(t-1)$, gives
$I(t)+C(t)\le C_3(I(t-1)+B(t))$ for $s\ge2$. Hence the weighted integral
is at least $cI(t)-CB(t)$.

\begin{lemma}[Two shrinking Q1 bands in the physical flow]\label{an04ce5v2:bandlemma}
Let $R$ be any sufficiently large fixed offset. Uniformly after a late fixed entrance,
\begin{equation}\label{an04ce5v2:uniform}
 \lim_{b\to\infty}\limsup_{q\downarrow0}\sup_{t_h+b\le t\le T} B(t)/I(t)=0 ,
 \qquad I(t_R-1)/q^{3/4}\longrightarrow\infty .
\end{equation}
\end{lemma}
\begin{proof}
We prove both assertions by testing the corresponding contrary sequences. First take $s=t-t_0\to\infty$ with $t\le t_R$ and $I(t)\le A B(t)$ for some fixed $A$. Use $\eps=(1+s)^{-4}$. In the second test take $t=t_R-1$, $I(t)\le A q^{3/4}$ and $\eps=q^{1/3}$. For either test write
\[
 g(u)=\sup_\alpha(|\mathsf R_{12}|+|\mathsf R_{21}|),\qquad
 P_t=\int_{t_0}^t g(u)\,du .
\]
By \eqref{an04ce5v2:prefix} and \eqref{an04ce5v2:earlyh}, $P_t=O_A(q(1+s))$ in the first test and $O_A(q^{3/4})$ in the second. In both,
\begin{equation}\label{an04ce5v2:units}
 q/\eps\to0,\quad P_t/\eps\to0,\quad
 \eps P_t/q\to0,\quad q e^{Qs/2}=O(1),\quad s=O(\ell).
\end{equation}
For the first test use $q(1+s)^5\to0$ and $(1+s)^{-3}\to0$; for the second the two perturbation ratios have powers $3/4-1/3$ and $1/3+3/4-1=1/12$. All comparisons below are simultaneous for $\eps\le k=\cos\alpha\le2\eps$ on Q1.

For coordinate $i$ set $X_i=(U_i,W_i)$, $J(x,y)=(y,x)$, and $a_i=\mathsf R_{ii}$. Exactly,
$X_i'=a_i JX_i+(\mathsf R_{j i}W_j,\mathsf R_{i j}U_j)$ with $j=3-i$.
Its homogeneous matrix is $e^{J\int a_i}$, of max-norm $e^{|\int a_i|}$. On equal positive entrance components it multiplies each by $e^{\int a_i}$. In sum/difference coordinates $(p,d)=((U_i+W_i)/2,(U_i-W_i)/2)$ the multipliers are $e^{\int a_i}$ and $e^{-\int a_i}$. We estimate the exact inhomogeneous integrals on whole prefixes; strict improvements on a closed trial region below exclude a first contact at a single face or a corner.

\paragraph{Initialization to the fixed entrance.}
Write $a_1^*=1/2+q^2,\ a_2^*=L+q^2$, $H=H_0,\ V=V_0$ as in Section~\ref{an04ce5v2:setup}; $V/H=e^{-\omega u}$. At $t_0$, $H\asymp_{b_0}1$ and $e^{-\omega t_0}\asymp_{b_0}q$ since $S_0=cq^{1/\omega}$. Up to $t_0$, $qH\le C_{b_0}V$, $\int H^2\le C_{b_0}$ and $\int V^2\le C_{b_0}q^2$. HBD's physical argument gives $|X_1|\le C H$, $|X_2|\le C(V+qH)\le C V$ and $g\le C(q+HV+qH^2)\le C(q+HV)$ in max-norm. For the fixed extension beyond HBD's cap $8$, one can use the same scalar inequalities there: the spectral bound $S\le H^2$, the individual radial bound integrated with $\int H^2\le C_{b_0}$, and $D^+\sup|X_2|\le(a_2^*+CH^2)\sup|X_2|+CqH(1+H^2)$ give the stated bounds by integration. Absolute fitted diagonals cost $C S$, since $|f(D)|\le S|D|$. The target diagonal is nonnegative and at most $a_i^*=\E D_i^2$, hence $|a_i|\le a_i^*+CS$. All constants in this paragraph may depend on $b_0$.

Variation of constants in horizontal units $H$ charges the perturbation on any prefix by at most
\[
 C\int_0^{t_0}(V/H)(q+HV)\,du\le C(q/\omega+\int V^2)=O(q).
\]
It follows first that $|X_1|\le C\eps H$ without a sign assumption. In vertical units $V$ the perturbation now costs at most
\[
 C\eps\int_0^{t_0}(H/V)(q+HV)\,du
 \le C\eps\{q e^{\omega t_0}/\omega+\int H^2\}=O(\eps).
\]
These bounds include the factor $\exp(C\int S)$ for both propagators.

For lower bounds subtract each missed target diagonal from $a_i^*$. For positive $U_i$ the missed diagonal on the principal cluster is bounded by $C\exp[-c(U_i/(q|U|))^2]$, by Cauchy--Schwarz, the Gaussian half-space probability, and the bounded fourth data moment. The other cluster costs at most $q^2/2$. Thus
$a_i\ge a_i^*-CS-q^2/2-C\exp[-c(U_i/(q|U|))^2]$.
Choose $M_0$ larger than the common bound for $C\int_0^{t_0}S+1$, before choosing any lower faces. Use the two lower faces for each layer $X_1\ge a\eps H$, $X_2\ge a V$, componentwise, where $0<a<e^{-M_0-2}/8$. At initialization $k\ge\eps$ and $\sin\alpha\ge1/2$ for small $q$, so all faces are strict. On the stopped region the upper estimates give strong margin $\ge c_a\eps/q$, weak margin
$\ge c_a\min(1/q,e^{\omega(t_0-u)}/\eps)$. The first integrated tail is $O(\ell e^{-c_a\eps^2/q^2})=o(1)$ for either band. The second is at most
\[
 C\ell e^{-c_a/q^2}
 +C\int_0^\infty e^{-c_a e^{2\omega v}/\eps^2}\,dv=o(1);
\]
use $e^{2\omega v}\ge1+2\omega v$. This bound also covers delays $s$ growing arbitrarily slowly. The $q^2t_0$ cost vanishes. Accordingly $e^{\int_0^u a_i}$ is at least $e^{-M_0-1}e^{a_i^*u}$ on every stopped prefix for small $q$. Equal entrance components give individual lower bounds in horizontal and vertical units of $\eps e^{-M_0-1}-O(q)$ and $e^{-M_0-1}/2-O(\eps)$. Both strictly improve their trial bounds for both components, including simultaneous contacts. Boundedness and continuity continue the comparison up to $t_0$.

\paragraph{From $t_0$ to the test time.}
Set $v=u-t_0$, $V_*(v)=q e^{Qv/2}$. At $v=0$ we have $\underline a\eps\le U_1,W_1\le\bar a\eps$ and $\underline a q\le U_2,W_2\le\bar a q$ with $\underline a>0$. Balance gives $|U_1-W_1|\le Cq^2/\eps=o(\eps)$ there. The finite-noise regression field gives
\[
 |a_1|\le C(|1-m|+q),\qquad
 a_2=\tfrac Q2(1-n)\theta_2+O(q),\qquad a_2\ge-Cq .
\]
Use the global initialized budgets $\int|1-m|+\int |n|\le C$, $1-n>0$ and $q s=o(1)$. For any subinterval the horizontal propagator has norm bounded by $C$, and the vertical one by $C e^{Q(v-v')/2}$. To justify both upper faces without assuming them, put $X=\sup |X_1|/\eps$, $Y=\sup |X_2|/V_*$ over a prefix and the band. The two integral equations give
\[
 X\le C+C\,(\sup V_*) P_t Y/\eps,\qquad
 Y\le C+C\,\eps P_t X/q .
\]
Their product of perturbation coefficients tends to zero by \eqref{an04ce5v2:units}. Both suprema are bounded, and the differences from the homogeneous solutions on every prefix cost at most $O(P_t)=o(\eps)$ in absolute horizontal units and $O(\eps P_t/q)=o(1)$ in relative vertical units. These constants are independent of the labels in the test band.

The homogeneous horizontal sum component is at least $c\eps$ and the absolute difference component is $o(\eps)$ even after propagation. This proves individual horizontal lower bounds. For the two vertical lower faces use a trial bound $U_2,W_2\ge a' V_*$, where $a'>0$ will be chosen below. On this region $U_2/(q|U|)\ge c_{a'}\min(1/q,e^{Qv/2}/\eps)$. Thus
\[
 \int_0^s(1-\theta_2)\,dv
 \le C s e^{-c_{a'}/q^2}
     +C\int_0^\infty e^{-c_{a'}e^{Qv}/\eps^2}\,dv=o(1)
\]
by the same weighted Gaussian tail bound. On any prefix write $A_2=\int_0^v a_2(t_0+v')dv'$. Then
$A_2\ge Qv/2-M_1-o(1)$, $A_2\ge-Cqs$, where $M_1$ bounds $(Q/2)\int|n|$ independently of $a'$. If the initial vertical sum/difference modes are $p_0,d_0$, $p_0-|d_0|\ge\underline a q$, $|d_0|\le Cq$. Both homogeneous components are at least
\[
 e^{A_2}(p_0-e^{-2A_2}|d_0|)
 \ge e^{A_2}(\underline a q-o(q)).
\]
Choose $a'<\underline a e^{-M_1-2}/4$ before taking $q$ small. After the inhomogeneous error $o(V_*)$ both components exceed $2a'V_*$, initially and at every possible first contact. The upper bounds were obtained unconditionally; the strict margins cover simultaneous lower contacts and all labels with either band endpoint. This proves the comparison through the full stopped horizon.

\paragraph{The physical overlap and the two contradictions.}
The Q1 band has label measure at least $\eps/(2\pi)$ since $|d(\cos\alpha)/d\alpha|\le1$. On the last unit, $U_1\gtrsim\eps$ and $U_2\gtrsim q e^{Qs/2}$; the unit factor $e^{-Q/2}$ is bounded below. Therefore $I(t)\ge c\eps^2 q e^{Qs/2}$. In the first test this contradicts $I(t)\le A q(1+s)$ since $e^{Ls}/(1+s)^9\to\infty$. In the second, $q e^{Qs/2}$ at $t_R-1$ is bounded below by a positive constant for fixed $R,b_0$, so $I(t)\gtrsim_R q^{2/3}$ contradicts $I(t)\le A q^{3/4}$. This proves the second assertion and rules out every contrary first-test sequence before $t_R$. Beyond $t_R$, monotonicity of $I$ bounds $B/I$ by $O(q(1+\ell))/I(t_R-1)=o(1)$. Failure of the displayed iterated upper limit would yield $b_j\to\infty$, $q_j\to0$ and times with $s_j\to\infty$ and $I=O(B)$, which the two ranges exclude. No resident limits on shrinking label sets enter these comparisons.
\end{proof}
Together with \eqref{an04ce5v2:prefix} and the two-unit estimate, this gives $\int_{t_0}^t z\ge c I(t)$ and $\int_{t_0}^t z_-/I(t)$ as small as desired uniformly after a sufficiently late fixed delay.


\section{Choosing a bulk without losing exposure to discarded labels}\label{an04ce5v2:retained}
\begin{lemma}[Physical weight of retained and terminally active labels]\label{an04ce5v2:bulklemma}
For sufficiently small fixed $\delta>0$, remove open axis strips of total label measure at most $\delta$ from $A_R$, leaving two compact interiors. There exist a further exceptional subset of measure $o_{q;\delta}(1)$ and a retained set $K$ in these interiors such that, on $[t_0,T]$, with $k_\alpha=\cos\alpha$,
\begin{equation}\label{an04ce5v2:bulk}
 D^{-1}k_\alpha\le U_1,W_1,b_1\le D k_\alpha,\qquad |U|+|W|\le Dk_\alpha
 \quad(\alpha\in K).
\end{equation}
Here $D$ and the coefficient $C$ in $\|\mathbf1_{A_R\setminus K}\mathcal U\|\le C\sqrt\delta+o_{q;\delta}(1)$ are independent of $\delta,R,b$. Moreover $\|\mathbf1_{A_R}\sup_t(|U_2|+|W_2|)\|=o_q(1)$ and $\|\mathbf1_{A_L\cap\{p>0\}}\mathcal U\|=o_q(1)$.
\end{lemma}
\begin{proof}
Throughout let $g$ be the uniform off-diagonal bound used in Lemma~\ref{an04ce5v2:bandlemma} (equivalently its defining supremum). First $G=\int_{t_0}^Tg=o_q(1)$. To check the limit use an \emph{auxiliary} cutoff $t_{R'}$ at large fixed $R'$. On $[t_0,t_{R'}]$, $g\le C(r+q)$; integration costs $Ce^{-LR'}+O(q\ell)$. On the remaining interval of length at most $R'+C$, $g\le C(|z|+|h|+q)=o_q(1)$ by Lemma~\ref{an04ce5v2:axislemma}. Take $q\to0$ at fixed $R'$, then let $R'\to\infty$. The same argument with $\|\mathbf1_{A_R}X_2\|$ in place of $g$ gives $\int_{t_0}^T\|\mathbf1_{A_R}X_2\|dt=o(1)$, using the incoming energy envelope first and uniform right axis separation second. Neither limit fixes the $R$ to be used for the gate comparison.

Since $|\mathsf R_{22}|\le C$ and $\|\mathcal U\|\le C$, integrate the two vertical equations labelwise and apply Minkowski after taking the time supremum:
\[
 \|\mathbf1_{A_R}\sup_t |X_2(t)|\|
 \le \|\mathbf1_{A_R}|X_2(t_0)|\|
     +C\int_{t_0}^T\|\mathbf1_{A_R} X_2\|dt+C G\|\mathcal U\|=o(1).
\]
The first term is $O_{b_0}(q)$. Max- and sum-norms of the two components differ by at most a factor two. For the horizontal pair write $A_{1,\alpha}(u,v)=\int_u^v\mathsf R_{11,\alpha}$, with $|A_{1,\alpha}|\le M$ on all subintervals by the integrated $|1-m|+q$ bound. Hence the horizontal difference from its exact homogeneous evolution has time supremum $\le C G\mathcal U_\alpha$, a function of vanishing label norm. On every compact right interior the entrance pair tends uniformly to $(a k_\alpha,a k_\alpha)$, $a=2\sqrt{M(b_0)}>0$, by HBD \texttt{clock} and BRW's fixed-delay tagged continuation. In particular, the maximum over $t$ of the error relative to the vector $e^{A_{1,\alpha}(t_0,t)}(a k_\alpha,a k_\alpha)$ has norm $o_{q;\delta}(1)$ there. The homogeneous multiplier is bounded between $e^{-M}$ and $e^M$ on this symmetric vector.

On the retained compact interiors discard labels where the preceding horizontal error or $\sup |X_2|$ exceeds $\epsilon_0 k_\alpha$, with $\epsilon_0<\min(1,a e^{-M}/4)$ fixed. By Chebyshev the discarded measure is $o_{q;\delta}(1)$ because $k_\alpha$ has a positive minimum on these \emph{fixed} compacts. On the good set both horizontal coordinates lie between $(a e^{-M}-\epsilon_0)k_\alpha$ and $(a e^M+\epsilon_0)k_\alpha$. The vertical upper bound and $|b_1-(U_1)_+|\le Cq|U|$ give all inequalities in \eqref{an04ce5v2:bulk} with one $D$ depending only on the global constants, $a,M,\epsilon_0$. No inverse minimum of $k_\alpha$ enters $D$. For any measurable $E\subset A_R$,
\[
 \|\mathbf1_E\mathcal U\|\le C\|\mathbf1_E X_1(t_0)\|+
 C G\|\mathcal U\|+C\|\mathbf1_{A_R}\sup_t |X_2(t)|\|
 \le C |E|^{1/2}+o_q(1)
\]
because all entrance horizontal coordinates have a common bound. Apply this to the axis strips and the additional discarded set. This proves small exceptional \emph{physical weight} in the norm of the time supremum, including its dependence on $\delta$.

On $A_L$, HBD's left mass bound gives $\|\mathbf1_{A_L}X_1(t_0)\|=O_{b_0}(q\ell)$. On any prefix the horizontal inhomogeneous solution costs at most $C\mathcal U_\alpha\int_{t_0}^t g$. Equations~\eqref{an04ce5v2:prefix} and \eqref{an04ce5v2:earlyh} therefore give
\begin{equation}\label{an04ce5v2:left}
 \|\mathbf1_{A_L}\sup_{[t_0,t]} |X_1|\|\le C(q\ell+I(t)).
\end{equation}
Over the full interval the same norm is $o(1)$ by $G=o(1)$. For the antisymmetric vertical mode $d=(U_2-W_2)/2$ the equation is $d'=-\mathsf R_{22}d+(\mathsf R_{12} W_1-\mathsf R_{21}U_1)/2$. Since $\mathsf R_{22}\ge-Cq$, its forward propagator is at most $e^{Cq(T-t_0)}$. Thus $\|\mathbf1_{A_L}\sup_t|d|\|\le C(q+G\|\mathcal U\|)=o(1)$. The excursion bound in Section~\ref{an04ce5v2:history} gives $\|\sup_t (U_2)_-\|=O(q\sqrt\ell)$. Hence $d_2=W_2-b_2$ also has vanishing left norm uniformly in time, using the feature error $Cq|U|$.

On $D_L=A_L\cap\{p>0\}$, $U_2(T)<\tan(q\zeta) U_1(T)\le Cq|U(T)|$. Combining this upper bound with the negative excursion gives $\|\mathbf1_{D_L}U_2(T)\|=o(1)$. Therefore the terminal norm there of $v=(U_2+W_2)/2=U_2-d$ is $o(1)$. Its coefficient is $+\mathsf R_{22}$ and its forcing $(\mathsf R_{12}W_1+\mathsf R_{21}U_1)/2$. The backward propagator is bounded by $e^{Cq(T-t_0)}$. Integrating backward labelwise and then taking norms gives $\|\mathbf1_{D_L}\sup_t|v|\|\le C(\|\mathbf1_{D_L}v(T)\|+G\|\mathcal U\|)=o(1)$. Together with the estimates for $d$ and $X_1$ this yields
\begin{equation}\label{an04ce5v2:activeleft}
 \|\mathbf1_{D_L}\mathcal U\|=o(1).
\end{equation}
Every comparison for a retained right label uses $k_\alpha$ as unit. Left comparisons concern the integrated physical weight; the bounds hold including the negative-horizontal initial directions.
\end{proof}


\section{Entry into the negative-input region}\label{an04ce5v2:entry}
Write $\nu=\inf_{[t_0,T]}(1-n)>0$ with a common lower bound for small
$q$. For $\alpha\in K$, \eqref{an04ce5v2:bulk} and the prefix estimates give,
after a late fixed delay $b$,
\[
 \int_{t_0}^t z b_1\,du\ge c_1 k_\alpha I(t)
 \qquad(t_h+b\le t\le T).
\]
Indeed the positive part is multiplied by at least $k_\alpha/D$ and
the negative part by at most $D k_\alpha$. The constant $c_1$ works
independently of $\delta$. We next pay for weak excitation up to
$t_b=t_h+b$. At each fixed $b$ the scaled weak-feature integral on a
retained right compact satisfies
\[
 \limsup_{q\downarrow0}\sup_K
 \frac{1}{q k_\alpha}\int_{t_0}^{t_b}|b_2|\,du
 \le C_\delta(1+b)e^{(L-\nu_\delta)b},\qquad \nu_\delta>0.
\]
This follows by $|b_2|\le (U_2)_++Cq|U|$, fixed-interval tagged
transfer, and the COD rank ceiling for $x_+$ on ranks bounded away from
the upper endpoint; the Q4 bottom path obeys the same bound. Take
$\nu_\delta<L$ if necessary. In contrast BRW's protected positive
first moment on the last unit gives
$\liminf_{q\to0} I(t_b)/q\ge c e^{Lb}/(1+b)$ for all large fixed $b$.
To transfer this lower bound, retain the positive overlap from the
right labels and use fixed-time extraction:
$(U_1)_+(U_2)_+/q$ converges there to the radial mass times $x_+$;
the nonnegative remainder may be discarded. TBW's bulk upper input
face also puts all retained compacts strictly below $U_2=0$ at $t_b$.
Thus, after fixing the compacts, choose $b$ sufficiently large for
the initial weak integral and $|W_2(t_0)|$ to cost less than any
prescribed fixed fraction of $k_\alpha I(t_b)$.

As long as $U_2\le0$ after this entrance on a retained label,
$|b_2|\le Cq k_\alpha$. The output equation integrates to
\begin{equation}
 W_2(t)\le -c_2 k_\alpha I(t),\qquad t_b\le t\le T,
 \label{an04ce5v2:wnegative}
\end{equation}
on this stopped interval: the additional weak cost is $Cq(t-t_b)k_\alpha$,
absorbed uniformly by \eqref{an04ce5v2:uniform}; the output shape error is
smaller. Explicitly, reserve $a_s=c_1/2$ from the lower bound for $(Q_1/2)\int z b_1/(k_\alpha I)$, since $Q_1\ge1$. Charge at most $a_s/8$ each to the entrance value of $W_2$, the weak term before $t_b$, the subsequent weak term, and the integrated shape error in these units. For the last two charges use respectively $CB/I$ and $CqB/I$ on each prefix. For the first two use the fixed-delay bounds above and $I(t)\ge I(t_b)$. The positive constant $c_2=a_s/2$ is fixed before $R,\delta,b$; only enforcement of the four fractions needs their subsequent choices. The input equation on that interval is
\begin{equation}
 U_2'=\tfrac Q2[(1-n)W_2-hW_1]\theta_2+O(qk_\alpha).
 \label{an04ce5v2:gate}
\end{equation}
Fix $\kappa>0$ sufficiently small compared with $\nu c_2/D$.
Whenever $h_-\le\kappa I$, \eqref{an04ce5v2:gate} gives
$U_2'\le-c_3 k_\alpha I\theta_2+Cq k_\alpha$.

For a sufficiently small fixed $d_g>0$ the stronger face needed at the final passage is
\begin{equation}\label{an04ce5v2:deepgate}
 U_2\le -d_g q k_\alpha\sqrt\ell .
\end{equation}
The next lemma proves entry by $t_R$ and continuation through $T$.

\section{Closing the final interval}\label{an04ce5v2:centered}
\begin{lemma}[Continuation of the negative input face]\label{an04ce5v2:gatelemma}
In the order of choices below, every label retained in Lemma~\ref{an04ce5v2:bulklemma} satisfies \eqref{an04ce5v2:deepgate} for $t_R\le t\le T$, and \eqref{an04ce5v2:wnegative} holds from $t_b$ through $T$.
\end{lemma}
\begin{proof}
The depth constant $d_g$ is distinct from the amplitude cutoff $d$ of Section~\ref{an04ce5v2:mass}. Fix $\nu>0$ from the passage bound $1-n\ge\nu$ and take $\kappa$ so that $\kappa D<\nu c_2/2$. Section~\ref{an04ce5v2:entry} gives the entry with these constants: on $U_2=0$ the principal weak gate is independent of $D_2$ and has probability $1/2$ (the retained $U_1>0$), so $\theta_2=1/2$. At $t_b$ the compact entrance lies strictly below zero. Up to $t_R$, \eqref{an04ce5v2:earlyh} gives $h_-/I\le C r(t_R)^2+CB/I<\kappa$ after $t_b$ with the choices below. While such an input is nonpositive, \eqref{an04ce5v2:wnegative} and this bound on $h_-$ give strict negative velocity at zero and on its interior side sufficiently near zero. One may also stop each individual label here because \eqref{an04ce5v2:earlyh} holds independently of other labels' gates. This gives $U_2\le0$ up to $t_R$. On $[t_R-1,t_R]$, at normalized input $u_\alpha=U_2/k_\alpha\in[-d_gq\sqrt\ell,0]$ the standardized weak normal mean has absolute value at most $\rho Dd_g\sqrt\ell$. For any chosen $\gamma\in(0,1/4)$ take $d_g$ so small that $(\rho Dd_g)^2/2<\gamma$. Integrating the standard normal density on an interval of length one beyond this threshold gives a lower gate probability $\ge c q^\gamma$. Removing $D_2\le\rho/2$ costs at most $Ce^{-c/q^2}$; on the remainder $D_2^2/Q\ge c$. Thus $\theta_2\ge c' q^\gamma$. The normalized speed is $\le-c I(t_R-1)q^\gamma+Cq$. Since $I(t_R-1)\gg q^{3/4}$, the magnitude of this inward speed exceeds the full normalized strip width divided by one unit of time for small $q$. Labels starting inside it cross the deep face before $t_R$; labels already below cannot cross upward. The inequalities include the endpoints and are uniform on $K$.

To continue, stop at the first time the supremum on $K$ of $u_\alpha$ reaches $-3d_gq\sqrt\ell/4$, or at $T$. Initially the supremum is at most $-d_gq\sqrt\ell$. It is Lipschitz in time because the field is bounded and $|W|/k_\alpha\le D$ on $K$. This gives a genuine stopped interval even without attaining the supremum. On it all retained labels are strictly negative. Their weak gate satisfies $|v_2|^2=o_q(1)$ uniformly: the margin magnitude is at least $3\rho d_g\sqrt\ell/(4D)$, Cauchy--Schwarz bounds the weighted principal tail by a constant times the square root of its probability, and $|(v_2)_1|\le Cq$. In particular the comparison for $W_2$ applies on every prefix from $t_b$.

Let $d_2=W_2-b_2$. The exact definition $T_2=\int W_1W_2|v_2|^2\dd$ and the feature bound $|b_2(1-|v_2|^2)|\le Cq|U|$ give $T_2-h=\int W_1|v_2|^2d_2\dd+O(q)$. Split this integral three ways. On $K$ it costs $o(1)\|W_1\|\|d_2\|$. On $A_R\setminus K$ it costs $(C\sqrt\delta+o_{q;\delta}(1))\|d_2\|$. On $A_L$ it costs $C(q\ell+I)\|\mathbf1_{A_L}d_2\|=o(1)(q\ell+I)$. Here $\|d_2\|\le C(I+B)$ by Section~\ref{an04ce5v2:prefixsec}. Since $I(t)\ge I(t_R-1)\gg q^{3/4}$ and $B=O(q\ell)$, uniformly on every stopped prefix,
\begin{equation}\label{an04ce5v2:t2}
 |T_2-h|\le(C_0\sqrt\delta+o_{q;R,\delta,b}(1))I(t).
\end{equation}
$C_0$ is independent of $R,\delta,b$; only the threshold for small $q$ can depend on all three.

Now the scalar CFCR equation with its $O(q)$ remainder gives
\[
 h'=a_h h+\tfrac12 Q(1-n)(T_2-h)+\tfrac12 Q_1(1-m)C(t)+O(q),
 \qquad a_h=\tfrac Q2(1-n-\|\nabla h\|^2).
\]
For every fixed $R$, $[t_R,T]$ lies in a fixed half-centered compact; \eqref{an04ce5v2:logistic} gives $n>0$ and $1-m=o_{q;R}(1)$ there. Consequently $a_h\le Q/2$. Also $C(t)\le C'(I+B)$, $q/I=o(1)$. By \eqref{an04ce5v2:earlyh},
$h_-(t_R)\le(C_1 e^{-2LR}+o_{q;R}(1))I(t_R)$, $C_1$ independent of $R$. Write $T-t_R\le R+C_T$ with $C_T$ fixed by the clock. The scalar propagator is positive. Its maximum is $e^{Q(R+C_T)/2}\le 2e^{L(R+C_T)}$ for small $q$ at fixed $R$. Applying variation of constants, using monotonicity of $I$ for both terms, gives on each stopped prefix
\[
 h_-(t)/I(t)\le C_2e^{-LR}+C_R\sqrt\delta+o_{q;R,\delta,b}(1),
 \qquad C_R=C_3(1+R+|C_T|)e^{L(R+|C_T|)}.
\]
The constants $C_2,C_3$ can be fixed independently of $R,\delta,b$. This estimate was derived assuming only the \emph{outer} stopped face; it improves the required $h_-$ bound there.

\begin{center}\small
\begin{tabular}{p{.14\linewidth}p{.76\linewidth}}
Choice & Dependencies and inequalities\\ \hline
First & Fix the family, $b_0$, passage cutoffs, $C_T,D,\nu,c_2,\kappa,\gamma,d_g$. The coefficient $c_2$ comes from the signed prefix with global $D,\nu$ and fixed relative error fractions.\\
$R$ & Take large for the incoming envelope, $Cr(t_R)^2<\kappa/4$ in \eqref{an04ce5v2:earlyh} and $C_2e^{-LR}<\kappa/4$. This fixes $C_R$.\\
$\delta$ & Keep a fixed central Q4 subarc and require $C_R\sqrt\delta<\kappa/4$, $C_5\sqrt\delta<c_4/2$ with the endpoint constants of Section~\ref{an04ce5v2:endpoint}.\\
$b$ & At fixed compacts, take $b$ large for TBW's negative entrance and the COD/BRW weak-cost comparison. Require the uniform $B/I$ and negative-$z$ costs small for the chosen fractions, including $CB/I<\kappa/4$.\\
$q$ & At fixed $R,\delta,b$ enforce strict tagged entrance, $t_b<t_R-1$, the band lower bound, and all remaining error fractions, in particular $|o_{q;R,\delta,b}(1)|<\kappa/4$.\\
\end{tabular}
\end{center}
The endpoint constants used here can be fixed before $\delta$: choose a closed central Q4 arc of positive length with $k_\alpha$ bounded below. Its retained measure is at least half its length for small $q$ at fixed truncation. Then $c_4$ depends only on this arc, $d_g,D,c_2$. $C_5$ uses the global $d_2$ and excursion bounds and the absolute physical strip constant of Lemma~\ref{an04ce5v2:bulklemma}. Neither requires an endpoint conclusion to set its value.

With these choices $h_-<\kappa I$ on the stopped interval. Equation~\eqref{an04ce5v2:gate} with $W_1\le Dk_\alpha$ gives $u_\alpha'\le-c I\theta_2+Cq$. Throughout the strip $-d_g q\sqrt\ell\le u_\alpha\le-3d_g q\sqrt\ell/4$ the same lower tail $\theta_2\ge c' q^\gamma$ applies, while the uniform upper-tail estimate already paid for \eqref{an04ce5v2:t2}. The derivative is strictly negative on the entire strip. The positive part $(u_\alpha+d_g q\sqrt\ell)_+$ starts at zero and cannot increase on the stopped interval, by its a.e. derivative, including both faces. Thus every retained label stays at or below the inner face. Taking the supremum leaves a positive distance to the outer stopping level even at a putative first exit. The stop must extend to $T$, including the case of exit and $T$ coinciding.
\end{proof}


\section{Subtracting all positive output outside the core}\label{an04ce5v2:endpoint}
On $K$, terminal exposure is at least $c q k_\alpha\sqrt\ell$, and
$-W_2(T)\ge c_2 k_\alpha I(T)$. The retained Q4 contains, for small
$\delta$ and small $q$, all but vanishing measure of a fixed central
Q4 subarc. Thus its favorable contribution is at least
$c_4 q\sqrt\ell\,I(T)$ with $c_4>0$ independent of the truncation.
Every retained Q1 label also has negative output.

To charge \emph{all} adverse labels, including any exceptional Q4,
let $D_*=\{p>0\}\setminus K$. Equations \eqref{an04ce5v2:activeleft} and
the right-strip estimate give
$\norm{\mathbf1_{D_*}\mathcal U}\le C\sqrt\delta+o_q(1)$.
The exposure excursion bound consequently gives
$\norm{\mathbf1_{D_*}p}\le(C\sqrt\delta+o_q(1))q\sqrt\ell$.
On the probe-active side $b_2\le Cq|U|$ because $U_2\le Cq|U|$.
Hence $\norm{\mathbf1_{\{p>0\}}(W_2)_+}\le\norm{d_2}+Cq\sqrt S
\le C I(T)$. Cauchy--Schwarz pays the entire adverse output by
$(C_5\sqrt\delta+o_q(1))q\sqrt\ell\,I(T)$.
Take $\delta$ small enough also that $C_5\sqrt\delta<c_4/2$.
The constants here precede the compact-specific choice of $b$.

Finally let $G_\alpha=\int_0^T[Q_1zb_1-w b_2]\dt$ and define
$M_q=(\int_{\rm Q4}pG_\alpha\dd-\int_{{\rm Q4}^c}p(-G_\alpha)_+\dd)/(2q^2)$.
The output projection error in replacing $-W_2(T)$ by $G_\alpha/2$
in these weighted integrals is $o(q^2)$ in absolute integral. Indeed
the shape error costs $O(q^2\ell)\norm p$, and the initial term costs
$O(\sqrt{S_0})\norm p$. Positive part is 1-Lipschitz. Consequently
\[
 M_q\ge c\sqrt\ell\, I(T)/q-o(1),\qquad
 H^2=O(\ell),\qquad I(T)/q^{3/4}\longrightarrow\infty .
\]
Since $\sqrt\ell\,I(T)/q\gg\ell$, this proves \eqref{an04ce5v2:criterion} with diverging margin. Directly, if
$F_i=f_{T,i}(\xi)$, then
$e_q=H^2/2-q^{-1}\sin(q\zeta)H+
(\cos(q\zeta)+F_2/2)F_2/q^2$.
The last coefficient tends to one by $\norm p=O(q\sqrt\ell)$.
Thus $e_q(T)\to-\infty$, which gives the strict comparison with
the finite coherent minimum at the measured learned endpoint. Since $F_i=O(q\sqrt\ell)$ the unscaled probe loss tends to $1/2$, as asserted in Theorem~\ref{an04ce5v2:main}.
\end{document}
